\documentclass[11pt,letterpaper]{article}
\usepackage[T1]{fontenc}
\usepackage[utf8]{inputenc}
\usepackage{lmodern}
\usepackage[margin=1in,headheight=15pt]{geometry}
\usepackage{amsmath,amssymb,amsthm,mathtools,mathrsfs}
\usepackage{microtype,booktabs,longtable,array,enumitem}
\usepackage{xcolor,listings,fancyhdr,needspace}
\definecolor{accent}{HTML}{4E6042}
\definecolor{ink}{HTML}{26312C}
\definecolor{codebg}{HTML}{F4F5F0}
\usepackage[colorlinks=true,linkcolor=accent,citecolor=accent,urlcolor=accent]{hyperref}
\usepackage[nameinlink,noabbrev]{cleveref}
\hypersetup{pdftitle={Foundations for Surreal and Surcomplex Mathematics},pdfauthor={Research exposition},pdfsubject={Set theory, type theory, size-safe analysis, and a Lean formalization architecture}}
\pagestyle{fancy}\fancyhf{}
\fancyhead[L]{\small Foundations for surreal and surcomplex mathematics}
\fancyhead[R]{\small\thepage}
\renewcommand{\headrulewidth}{0.3pt}
\setlength{\parskip}{0.28em}
\setlength{\emergencystretch}{3em}
\setlist{itemsep=0.15em,topsep=0.35em,leftmargin=1.7em}
\setcounter{tocdepth}{2}
\numberwithin{equation}{section}
\newtheorem{theorem}{Theorem}[section]
\newtheorem{lemma}[theorem]{Lemma}
\newtheorem{proposition}[theorem]{Proposition}
\newtheorem{corollary}[theorem]{Corollary}
\theoremstyle{definition}
\newtheorem{definition}[theorem]{Definition}
\newtheorem{example}[theorem]{Example}
\newtheorem{principle}[theorem]{Design principle}
\theoremstyle{remark}
\newtheorem{remark}[theorem]{Remark}
\newcommand{\No}{\mathbf{No}}
\newcommand{\SC}{\mathbf{No}[i]}
\newcommand{\On}{\mathbf{On}}
\newcommand{\Oz}{\mathbf{Oz}}
\newcommand{\R}{\mathbb{R}}
\newcommand{\C}{\mathbb{C}}
\newcommand{\Q}{\mathbb{Q}}
\newcommand{\Z}{\mathbb{Z}}
\newcommand{\N}{\mathbb{N}}
\newcommand{\OO}{\mathcal{O}}
\newcommand{\mm}{\mathfrak{m}}
\newcommand{\UU}{\mathcal{U}}
\newcommand{\supp}{\operatorname{supp}}
\newcommand{\st}{\operatorname{st}}
\newcommand{\cis}{\operatorname{cis}}
\newcommand{\fin}{\operatorname{fin}}
\newcommand{\rk}{\operatorname{rank}}
\newcommand{\bd}{\operatorname{b}}
\newcommand{\hSum}{\mathop{\sum}\nolimits^{\!H}}
\newcommand{\Class}{\mathsf{Class}}
\newcommand{\Small}{\mathsf{Small}}
\newcommand{\Type}{\mathsf{Type}}
\newcommand{\Prop}{\mathsf{Prop}}
\newcommand{\ZF}{\mathsf{ZF}}
\newcommand{\ZFC}{\mathsf{ZFC}}
\newcommand{\GB}{\mathsf{GB}}
\newcommand{\GBC}{\mathsf{GBC}}
\newcommand{\KM}{\mathsf{KM}}
\newcommand{\ETR}{\mathsf{ETR}}
\newcommand{\Con}{\operatorname{Con}}
\newcolumntype{L}[1]{>{\raggedright\arraybackslash}p{#1}}
\lstset{basicstyle=\ttfamily\footnotesize,columns=fullflexible,keepspaces=true,
  breaklines=true,breakatwhitespace=false,showstringspaces=false,
  backgroundcolor=\color{codebg},frame=single,rulecolor=\color{black!15},
  xleftmargin=0.5em,xrightmargin=0.5em,aboveskip=0.7em,belowskip=0.7em,
  literate={→}{{$\to$}}1 {∀}{{$\forall$}}1 {∃}{{$\exists$}}1 {¬}{{$\neg$}}1
  {∈}{{$\in$}}1 {ℝ}{{$\R$}}1 {ℚ}{{$\Q$}}1 {ℕ}{{$\N$}}1
  {≤}{{$\le$}}1 {≠}{{$\ne$}}1 {⟨}{{$\langle$}}1 {⟩}{{$\rangle$}}1
  {↔}{{$\leftrightarrow$}}1 {×}{{$\times$}}1 {·}{{$\cdot$}}1
  {²}{{$^2$}}1 {Γ}{{$\Gamma$}}1 {α}{{$\alpha$}}1 {β}{{$\beta$}}1}
\begin{document}
\begin{titlepage}
\centering
\vspace*{0.8cm}
{\large\scshape Foundations and formal verification\par}
\vspace{0.8cm}
{\Huge\bfseries\color{ink} Foundations for Surreal\\[0.18em] and Surcomplex\\[0.18em] Mathematics\par}
\vspace{0.8cm}
{\Large Set Theory, Type Theory,\\[0.2em] Size-Safe Analysis, and Lean\par}
\vspace{0.7cm}
{\large September 21, 2026\par}
\vspace{0.7cm}
\begin{minipage}{0.93\textwidth}
\begin{abstract}
Surreal numbers are not contradictory large real numbers. They are a precisely constructible ordered field whose totality is too large to be a set in its own ambient set universe. The real foundational danger is to erase this qualification: unrestricted cut formation, naive class quotients, proper-class supports, and universe-insensitive topology lead to false or inconsistent statements. We compare definable-class constructions in ZF and ZFC, two-sorted class theories, set-sized fragments and Hahn fields, Grothendieck universes, constructive set theories, dependent type theory, and higher inductive-inductive constructions. We distinguish the existence of the underlying field from axiomatizations that characterize its simplicity hierarchy and from additional analytic structure.

The supplied analysis and trigonometry manuscripts serve as concrete test cases. We give size-correct formulations of their workspace, summability, topology, and phase-extension principles, including a precise explanation of why small-index nets may be eventually constant while larger-index nets need not be. For Lean, we inspect a pinned contemporary combinatorial-games revision, identify its existing surreal field and small-support Hahn infrastructure, and propose a layered implementation that separates generic algebra, support-controlled analysis, and the Conway normal-form bridge. Illustrative Lean files and a source-audit record accompany the article. They are not represented as an executed formal verification.
\end{abstract}
\end{minipage}
\vfill
\begin{minipage}{0.93\textwidth}\small
\textbf{Meaning of ``paradox-free.''} All consistency assurances are relative to the chosen foundational theory and its metatheory. This article supplies constructions, explicit size restrictions, and proof obligations; it does not claim an absolute proof of the consistency of set theory or Lean.
\end{minipage}
\end{titlepage}
\begingroup\small\setlength{\parskip}{0pt}\tableofcontents\endgroup
\clearpage

\section{Scope, sources, and the central recommendation}\label{sec:scope}
The problem has three layers that should not be merged. The first is a \emph{foundational construction}: what objects are surreal numbers, what does it mean to quantify over them, and why are recursive definitions legitimate? The second is a \emph{mathematical interface}: which properties of ordered fields, cuts, simplicity, normal forms, and exponentiation are available? The third is \emph{formal verification}: how can a proof assistant represent these objects and check arguments without silently changing their size or their topology?

The classical underlying facts are that $\No$ is a real-closed ordered class field and that $\SC$ is an algebraically closed class field \cite{Conway,Gonshor,Alling}. The adjective \emph{class} qualifies the carrier, not the arity of addition, multiplication, or polynomial equations. Each individual number has a set-sized description. Each polynomial has an ordinary finite degree. There is no need for proper-class tuples of coefficients to formulate algebraic closedness.

\subsection{What the two supplied manuscripts actually assume}
The supplied \emph{Surcomplex Analysis} manuscript explicitly chooses a class theory such as NBG with global choice, requires all normal-form supports and all coefficient families to be sets, and uses set-sized Hahn workspaces. It singles out arbitrarily small rescaling and all-scale rigidity as places where a previously chosen workspace may be insufficient \cite[\S2.1 and \S4.1]{AnalysisSource}. Its fine topology is described by balls with arbitrary positive surreal radii, not by a purported set of all open subclasses.

The supplied \emph{Trigonometry on the Surcomplex Plane} manuscript retains the same support convention. Its circular functions are initially defined on finite surreal arguments, its geometric lengths range over all surreals, and its extensions to infinite circular arguments require extra phase data \cite[\S\S2--4 and 16]{TrigSource}. These are substantive domain restrictions, not presentational conveniences.

Our discussion is an expansion and foundational audit of those conventions, not a claim to have machine-checked the manuscripts. Statements attributed to them are identified as such. Published results, conclusions proved here, and proposed implementation interfaces are kept separate. The other archives mentioned in the trigonometry bibliography were not part of the materials examined for this article.

\subsection{A practical answer before the details}
For a conventional mathematical exposition, a carefully specified version of GBC is a convenient language: individual numbers and supports are sets; $\No$, its operations, and particular global analytic functions are classes. The basic theory can instead be expressed in ZFC through definable classes, without a new large-cardinal axiom merely for surreal arithmetic.

For Lean, the most effective architecture is usually different. Work with a universe-indexed surreal type, prove theorems about \emph{small} families at that level, and implement the analytic core over ordinary set-sized Hahn fields. Reuse the current game-based surreal field rather than starting its algebra from zero. Prove explicit comparison theorems between the game, sign, and Hahn presentations. The full normal-form theorem and analytic enrichments should be tracked as mathematical dependencies, not inferred from the existence of suggestively named modules.

\section{The size obstructions that every foundation must respect}\label{sec:obstructions}
\subsection{The impossible unrestricted cut axiom}
Write $L<R$ for $\forall l\in L\,\forall r\in R\;(l<r)$. The familiar surreal rule constructs a number $\{L\mid R\}$ between \emph{set-sized} separated families. Omitting that adjective changes the rule fundamentally.

\begin{theorem}[No carrier is closed under all of its own cuts]\label{thm:nocut}
Let $F$ be a collection with an irreflexive relation $<$. If every subclass $L\subseteq F$ is an admissible left family and the empty right family is admissible, the assertion
\[
 \forall L\subseteq F\ \exists x\in F\ \forall l\in L\;(l<x)
\]
is inconsistent. In particular, no set can satisfy the surreal cut axiom for all of its subsets.
\end{theorem}
\begin{proof}
Take $L=F$. A proposed output $x\in F$ must satisfy $x<x$. This contradicts irreflexivity. The same reasoning works internally in type theory whenever the rule accepts the entire carrier as an indexing family.
\end{proof}

The legitimate class-field rule cannot be applied to $L=\No$, because $\No$ is not a set. Likewise, a universe-indexed constructor that accepts only $u$-small families cannot accept the full type $\No_u$, which lies at a higher level and is not $u$-small. The restriction blocks the proof exactly at the invalid input; it is not an exception added after obtaining a contradiction.

\subsection{Four different assertions involving ``all''}
These statements have different meanings:
\begin{enumerate}
\item Every ordinal of the ambient set universe embeds in the surreal class.
\item Every ordinal small relative to a specified universe embeds in its associated surreal type.
\item Every set of surreal coefficients fits in a sufficiently large set-sized workspace.
\item There is one set-sized workspace containing every surreal number.
\end{enumerate}
The first three have appropriate rigorous formulations. The fourth contradicts the first. In particular, a theorem producing a workspace \emph{after} a set of inputs has been given does not produce a universal workspace.

\begin{proposition}[Properness from birthdays]\label{prop:proper}
In the sign-sequence construction, the class of all surreal numbers is not a set.
\end{proposition}
\begin{proof}
For each ordinal $\alpha$, the all-plus sign sequence of length $\alpha$ is a surreal with birthday $\alpha$. If the total carrier were a set, Replacement applied to the birthday map would produce a set containing every ordinal. No such set exists: the union of a set of ordinals is an ordinal, and its successor is larger than every ordinal in that set after taking a suitable upper bound.
\end{proof}

\subsection{A universal set is not a shortcut}
Foundations with a universal set, such as stratified set theories, do not license the unrestricted rule in \cref{thm:nocut}. They restrict comprehension, type assignments, or the class of admissible constructions. A theory with a universal carrier and a strict order cannot also prove that every subset of that carrier has a strict upper bound inside it. Thus a change of foundational language must explain \emph{which} input families remain admissible. It cannot remove this elementary obstruction.

Nor is $\No$ a Dedekind-complete ordered field. Let $H>n$ for every $n\in\N$. Then $\N$ is bounded above. If it had a least upper bound $s$, the smaller number $s-1$ would still be an upper bound, because $s\ge n+1$ for every $n$. Surreal cut interpolation is therefore not least-upper-bound completeness. The distinction matters both in mathematical prose and when selecting Lean typeclasses.

\section{ZF and ZFC: constructing a proper class without class objects}\label{sec:zfc}
\subsection{Sign sequences as individual set codes}
A particularly transparent carrier is the definable class
\begin{equation}\label{eq:signclass}
 \No=\{s:\operatorname{dom}(s)\text{ is an ordinal and }s:\operatorname{dom}(s)\to\{-,+\}\}.
\end{equation}
Each $s$ is an ordinary set-theoretic function. Its birthday is $\bd(s)=\operatorname{dom}(s)$. To compare two distinct sequences, take the first place at which they differ, treating termination as a third symbol ordered by
\[
 -\ <\ \text{termination}\ <\ +.
\]
Initial segment extension is the simplicity relation, not the numerical order. Both the ordering and the tree structure are definable using set quantifiers. Gonshor's presentation is a standard reference for this representation \cite{Gonshor}.

For any fixed ordinal $\alpha$, the collection
\[
 \No_{<\alpha}=\bigcup_{\beta<\alpha}\{-,+\}^{\beta}
\]
is a set. Every set $A$ of sign sequences is contained in some such stage: take an ordinal strictly above the supremum of its birthdays. This simple Replacement argument is the source of many later localization arguments.

The ordinal embedding sends $\alpha$ to its all-plus sequence. It is an order embedding, but it does not identify ordinal addition with surreal field addition. For example, surreal addition is commutative and $\omega+1>\omega$, while ordinal $1+\omega=\omega$. Formal libraries must use separate operations and explicit comparison results rather than treating an ordinal coercion as a homomorphism for ordinary ordinal arithmetic.

\subsection{What ``class'' means in ZFC}
ZFC has only set variables. In that language, $x\in\No$ abbreviates a formula $\operatorname{Surreal}(x)$. A class function $f:\No\to\No$ is represented by a formula $\varphi_f(x,y)$ for which one proves
\[
 \forall x\bigl(\operatorname{Surreal}(x)\Rightarrow
   \exists!y\,[\operatorname{Surreal}(y)\land\varphi_f(x,y)]\bigr).
\]
There is no set of all inputs or of all graph pairs. Nevertheless, each restriction to a set of inputs has a set of outputs by Replacement. A finite formula expressing a field law can quantify over surreal codes without quantifying over a new kind of object.

This route is enough for definitions of the canonical arithmetic operations and for set-parameterized theorems about them. It is less convenient for a theorem that quantifies over \emph{arbitrary class functions}: such a statement becomes a schema over definable formulas, or requires an expanded language. Confusing a schema about definable functions with quantification over all classes can strengthen or weaken a statement unintentionally.

\subsection{Recursion: numerical size is not the termination measure}
For a sign sequence $x$, its canonical left and right options are suitable shorter initial segments. All options are simpler than $x$, although right options are numerically larger. Recursive addition has the familiar form
\begin{equation}\label{eq:conwayadd}
 x+y=\{x^L+y,\ x+y^L\mid x^R+y,\ x+y^R\}.
\end{equation}
A well-founded relation on pairs of birthdays, or on pairs of option trees, justifies the recursion. The relation $<$ on numerical values cannot do so: $0>-1>-2>\cdots$ is an infinite descending chain.

For each input pair, the recursively consulted subpositions form a set. One first solves the recursion on a set-sized, downward-closed domain containing those subpositions, proves that any two such local solutions agree on their overlap, and then uses the unique local output as the class definition. Multiplication requires additional algebraic expressions in earlier values, but its recursive references still decrease an input's simplicity. Inversion is more involved and needs its own well-founded construction; it does not follow merely from writing a suggestive cut.

A rigorous development should record three independent claims at each recursive definition: the recursive calls are well founded, the option families are sets, and the constructed left options lie below the right options. Well-foundedness alone does not imply that a game is a number.

\subsection{What is and is not being claimed about weak axioms}
The sign-sequence carrier and its order can be described without the axiom of choice. This does not establish that every analytic theorem in the supplied manuscripts has a choice-free proof. Here the complete classical working theory is taken to be ZFC, with choice dependencies isolated when relevant. Identifying the weakest subsystem for the entire development would be a separate reverse-mathematical investigation.

Replacement is not decorative. It bounds the birthdays and collects the supports of a set family. Weakening the set theory requires rechecking these steps. Even removing the power-set axiom has subtleties: Replacement and Collection cease to have some of their familiar equivalences \cite{GHJ}. No claim that arbitrary weak fragments support the full Hahn and analytic theory is intended.

\section{Conway names, quotients, and canonical representatives}\label{sec:quotients}
\subsection{Raw names are not yet numbers}
A raw Conway name is a well-founded tree with a set of left children and a set of right children. The numeric condition says recursively that the children are numeric and every left value is less than every right value. Distinct names can represent the same number. For example,
\[
 \{\varnothing\mid\varnothing\}\quad\text{and}\quad\{-1\mid1\}
\]
both represent zero, although their tree structures and birthdays as presentations differ.

The equivalence relation on numeric names is numerical equality:
\[
 x\approx y\quad\Longleftrightarrow\quad x\le y\ \land\ y\le x.
\]
One must prove that the operations respect this relation. Only then may they descend to numbers. Equality of raw trees is not a replacement for this congruence proof.

\subsection{The class-quotient trap}
In a set theory, quotienting a \emph{set} by an equivalence relation produces a set of equivalence classes. For a proper class of raw names, individual equivalence classes need not be sets. Zero alone has numeric presentations of unbounded rank: for arbitrarily large positive ordinals $\alpha$, the name $\{-\alpha\mid\alpha\}$ represents zero. Thus the notation ``the class of equivalence classes'' cannot simply mean a class whose elements are those proper classes. In GB or NBG, proper classes are not elements.

There are several sound repairs. Use canonical sign sequences as the final carrier. Alternatively, normalize each raw name to a canonical code and quotient by equality of those codes. A more general device is Scott's trick, which avoids selecting a single representative.

\begin{proposition}[Scott codes for a definable class quotient]\label{prop:scott}
Let $E$ be a definable equivalence relation on a definable class $D$ of sets. For $x\in D$, let $\rho_x$ be the least rank of a set $E$-equivalent to $x$, and define
\[
 q(x)=\{y\in D:\rk(y)=\rho_x\text{ and }yEx\}.
\]
Then $q(x)$ is a nonempty set and $q(x)=q(z)$ exactly when $xEz$.
\end{proposition}
\begin{proof}
The rank of $x$ supplies an upper bound for the search for $\rho_x$, so the least rank exists. All candidates of that rank lie in $V_{\rho_x+1}$, and Separation forms $q(x)$. Equivalent inputs have exactly the same candidate class and hence the same code. Conversely, a member of a common nonempty code is equivalent to both inputs, which implies $xEz$.
\end{proof}

Scott coding preserves an extensional quotient without a global choice of representatives. It is not necessarily the most efficient representation for arithmetic: transporting recursive functions and simplification through these codes adds overhead. This is an engineering reason to prefer signs or a native type-theoretic quotient, not a foundational defect in the construction.

\subsection{Canonical normalization and choice}
A uniquely specified simplest representative may be definable and require no selection axiom. An arbitrary representative from each equivalence class is a different matter. In a class theory, making arbitrarily many such selections simultaneously can require a global choice principle. In a fixed Lean universe, the quotient and its fibers are types, and \texttt{Quotient.out} uses Lean's choice axiom at that level. The logical roles resemble one another, but ``Lean choice'' and ``NBG global choice'' are not literally the same axiom in the same language.

The Mizar development of P\k{a}k and Kaliszyk demonstrates a substantial alternative proof-engineering route: it connects presentations, replaces an induction-recursion organization with ordinary transfinite induction, and develops field operations and Conway normal form \cite{PakKaliszyk}. Its existence is evidence that a proof assistant need not have a primitive surreal inductive-inductive constructor to formalize the classical theory.

\section{Class theories: GB, GBC, NBG, and Kelley--Morse}\label{sec:classes}
\subsection{Fix the convention before invoking a theory name}
The names GB and NBG are used with slightly different conventions concerning choice. In this article, \emph{GB} means a two-sorted G\"odel--Bernays theory with sets and classes, elementary class comprehension, and the usual class-replacement principle; the choice assumptions are stated separately. \emph{GBC} adds global choice and has ZFC as its set part. Many uses of ``NBG with global choice'' describe the corresponding classical working framework. None of our comparisons depends on treating the letters NBG as a guarantee of a particular unstated choice convention.

A class is a collection of sets. Only sets can be members. Elementary class comprehension allows class parameters, but its defining formula quantifies only over sets. Thus, for a given class function $f$, the class of points where $f$ satisfies a specified algebraic or epsilon--delta condition can be formed without impredicative comprehension, provided all its bound variables are sets. Quantifying over surreal numbers is set quantification, because each surreal is a set code.

GBC is conservative over ZFC for sentences in the language of sets. This does not make the theories identical: GBC can quantify over classes, and models can have different collections of classes over the same underlying set universe. The literature on second-order set theories makes this distinction explicit \cite{Williams}. Conservativity also does not say that every particular model of ZFC already has a definable global well-order.

\subsection{Why the manuscripts fit naturally in a class language}
The graphs of addition and multiplication are classes of finite tuples of sets. The finite elements $\OO\subset\No$, the infinitesimals $\mm\subset\No$, the surcomplex unit circle, and the purely infinite part $\Pi$ are classes. A specific exponential, derivative relation, or phase character is also a class of appropriate set-coded tuples. These objects need not be assembled into a set in order to state their laws.

The assertion ``for every class homomorphism $\chi:\Pi\to\mathbb T(\No)$'' is a legitimate second-sort quantification. It is not the assertion that there exists a class \emph{whose elements are all such homomorphisms}. The latter would require classes to be elements, hence a higher-order framework or a coding restriction. This is the proper reading of the phase classification in the trigonometry manuscript \cite[\S16.1]{TrigSource}.

Similarly, a topology on a proper-class carrier can be discussed through a neighborhood relation or through a formula saying that a specified class is open. There is generally no object $\mathcal P(\No)$ collecting all subclasses, and no ordinary set $\tau\subseteq\mathcal P(\No)$ of all opens. On a set-sized workspace, of course, an ordinary topology is a set and the usual definition is available.

\subsection{Two fundamentally different kinds of recursion}
\begin{principle}[Audit the output size as well as the index size]\label{prin:recursion}
A recursion with a proper-class domain can still be a local, set-valued recursion. Conversely, a recursion of ordinal length can ask for a new \emph{class} at each stage and require a substantially stronger principle. The length of the index alone does not determine the strength needed.
\end{principle}

Suppose a recursive rule computes a set $F(x)$ from values at predecessors of $x$, and each downward dependency domain is a set. A rank argument proves existence and uniqueness of the local solution. Its graph is a class obtained by elementary comprehension from the assertion that the unique local solution has a certain value. This is the pattern behind ordinary well-founded recursive constructions of surreal operations. A set-like well-founded relation has only set-many predecessors per element, and is the standard favorable case \cite{HamkinsRecursion}.

By contrast, elementary transfinite recursion for \emph{class-valued} stages includes iterations of truth predicates. At a stage one may need a relation on all set assignments, rather than one set. GBC does not in general prove the full ETR principle; stronger class theories can. Even an Ord-indexed construction can exhibit this difficulty when the stages are classes. Results on class forcing locate important recursions at $\ETR_{\On}$ \cite{ClassForcing}.

It is therefore wrong to justify every construction with the sentence ``NBG permits transfinite recursion.'' It is equally wrong to require Kelley--Morse solely because $\No$ is a proper class. One must identify the actual recursion and prove that its dependency domains and outputs have the advertised sizes.

\subsection{What Kelley--Morse adds}
Kelley--Morse allows impredicative class comprehension: class quantifiers may occur in the formula defining a class. This supports stronger class-recursion and satisfaction constructions. It is not merely a conservative notational extension of ZFC comparable to GBC; its proof-theoretic and model-theoretic behavior is different \cite{Williams,ClassForcing}.

KM can be a legitimate ambient theory, especially for a project about classes of structures, iterated truth, or sophisticated class recursions. But adopting it does not solve the universal-cut contradiction or permit proper classes to become ordinary members. If a proposed argument requires a hyperclass of all class functions, that size issue still has to be addressed. Stronger comprehension is not unrestricted type collapse.

\subsection{A choice audit for the present subject}
The simplest element of a separated surreal cut is unique, so its selection is canonical. Taking the union of a set of supports also needs no choice of representatives once normal forms are defined. Selecting roots of a set of polynomials can use ordinary choice, although many individual roots may instead be canonically selected by ordering or additional data. Choosing representatives or bases across a proper class of unrelated structures can require global choice.

None of this entails that every proof using global choice essentially needs it. For the explicit phase examples in the trigonometry manuscript, a coefficient functional on normal forms gives the character directly; no proper-class basis is required. An abstract proof using a basis of a class-sized vector space would need a separate justification and should not be substituted without notice.

\section{Set-sized fragments and universe enlargement}\label{sec:fragments}
\subsection{Birthday truncation is not automatically algebraic closure}
The stages $\No_{<\alpha}$ are excellent for construction and induction, but an arbitrary stage is not a field. For example, the surreals born at finite days are dyadic rationals. They contain $3$ but not $1/3$. Larger bounds require actual closure theorems for addition, multiplication, inversion, roots, and exponentiation; the needed bounds need not coincide.

A formalization should not insert a field instance on a birthday subtype merely because its ambient type is a field. It must prove closure of the subtype under every operation in that instance. Results specifying suitable ordinal bounds exist in the surreal literature, including work on exponential subfields \cite{vdDE,EKExp}; these are mathematical theorems, not consequences of notation. Exact bounds should be ported together with their hypotheses and relevant corrections.

\begin{proposition}[Closure under a fixed set of finitary operations]\label{prop:closure}
Let $A$ be a set of surreal codes. Let $\mathcal F$ be a uniformly specified set-indexed family of finitary class operations, with set values on their domains. More precisely, one class graph encodes the operations together with their set indices and finite arities; the operations themselves are not members of a set. Then a set $B\supseteq A$ exists that is closed under every operation in $\mathcal F$ whenever the inputs belong to $B$.
\end{proposition}
\begin{proof}
Put $A_0=A$. Given $A_n$, adjoin every value of every operation in $\mathcal F$ on tuples from $A_n$ lying in its domain, obtaining $A_{n+1}$. Finite Cartesian powers, Replacement, and Union ensure that this is a set. Let $B=\bigcup_{n<\omega}A_n$. Every finite tuple from $B$ already lies in one $A_n$, so its output lies in $A_{n+1}\subseteq B$.
\end{proof}

For field operations, include $0,1,+,-,\cdot$ and inversion away from zero. For canonical exponential and logarithmic closure, include those particular operations. This proposition does \emph{not} produce closure under all set-indexed Hahn sums: those operations have unbounded arity and their admissible input families are not one fixed finitary signature. A full Hahn field is a more suitable workspace when such sums are central.

\subsection{Grothendieck universes and inaccessible cardinals}
Suppose the ambient metatheory contains a Grothendieck universe $\UU$, for example $V_\kappa$ for a strongly inaccessible cardinal $\kappa$. Inside it, one can construct all surreal sign sequences whose ordinal lengths and graphs belong to $\UU$. From the larger universe these constitute a set $\No^{\UU}$, typically not an element of $\UU$ itself. It contains the ordinals of $\UU$, not every ordinal of the external metatheory.

If $A\in\UU$ is a family of its elements, the birthdays of $A$ are bounded below $\kappa$, and the corresponding cut can remain in $\No^{\UU}$. But the external set of \emph{all} members of $\No^{\UU}$ is not an admissible $\UU$-small family. Thus \cref{thm:nocut} is avoided in exactly the same way as before.

A set-theoretic universe assumption has a cost. The existence of a strongly inaccessible cardinal is not a theorem of ZFC. Requiring arbitrarily large such universes is stronger still. These assumptions may be convenient for interpreting a hierarchy of type universes, but they are not needed merely to discuss the ordinary definable proper class $\No$ in ZFC. A semantic model of a type theory is also not a theorem that its exact proof-theoretic strength equals that of one familiar large-cardinal axiom.

\subsection{Tarski--Grothendieck and Mizar}
Tarski--Grothendieck set theory provides a universe-rich environment extending ZFC. This is the foundation used by the Mizar formalization cited above \cite{PakKaliszyk}. A completed Mizar proof establishes a result in that formal environment. Extracting a proof in a weaker theory requires a dependency analysis; it does not follow simply because the informal theorem is usually stated in ZFC.

For a Lean project whose aim is surreal geometry rather than foundational comparison, implementing set theory inside type theory solely to reproduce a TG-style construction adds a large translation burden. For a project whose aim is to compare theories, however, that additional layer is precisely the object of study and can be appropriate.

\subsection{External countability and internal properness}
A countable model of set theory can internally believe that its surreal collection is a proper class. Externally, its domain and its internal surreal codes may be countable sets. There is no contradiction: the external enumeration need not be an internal set or class function available in the model. Similarly, ``every set family has a bound'' means every family represented by a set \emph{in the relevant model}. It is not an absoluteness claim about every family visible from a stronger metatheory.

This warning applies to normal forms and supports as well. Set membership, well-foundedness, completeness, and the collection of all reals must be interpreted in the stated ambient model. Statements about standard mathematical surreals usually use the intended external universe; model-theoretic comparisons must make the interpretation explicit.

\section{Hahn fields: a modular set-sized analytic foundation}\label{sec:hahn}
\subsection{The workspace construction}
Let $\Gamma$ be a set-sized linearly ordered abelian group. A Hahn series over a field $k$ is a function $a:\Gamma\to k$ whose nonzero support is well ordered. We write
\[
 k((t^\Gamma))=\left\{\sum_{\gamma\in S}a_\gamma t^\gamma:
                S\subseteq\Gamma\text{ well ordered}\right\}.
\]
This is a set, constructed as a subclass of the function set $k^\Gamma$. Addition is coefficientwise. Multiplication is convolution; the support theorem ensures that each coefficient is a finite sum. The valuation is the least exponent of a nonzero series.

With the convention of the supplied manuscripts,
\[
 t^\gamma=\omega^{-\gamma},\qquad
 F_\Gamma=\R((t^\Gamma)),\qquad K_\Gamma=\C((t^\Gamma)).
\]
For a divisible $\Gamma$, $K_\Gamma$ is algebraically closed by the standard Hahn-field theorem \cite[Corollary 4]{Poonen}. Splitting real and imaginary coefficients identifies $K_\Gamma$ with $F_\Gamma[i]$. Since $F_\Gamma$ is ordered and its quadratic extension by $i$ is algebraically closed, $F_\Gamma$ is real closed. These are set-sized fields to which ordinary algebraic formalization tools apply.

\begin{theorem}[Capturing any prescribed set of surreal data]\label{thm:workspace}
Assume the Conway normal-form theorem. Every set $A\subseteq\SC$ is contained in a canonically embedded field $K_\Gamma$ for some set-sized divisible ordered subgroup $\Gamma\subseteq(\No,+)$.
\end{theorem}
\begin{proof}
Collect the real and imaginary supports of all members of $A$, with their signs reversed to match $t^\gamma=\omega^{-\gamma}$. Their union $S$ is a set by Replacement and Union. Let $\Gamma$ be the $\Q$-linear span of $S$ in $\No$; it consists of finite rational linear combinations and is therefore a set. It is divisible and inherits its order from $\No$. Every member of $A$ has support in $\Gamma$. Evaluating a Hahn series by its Conway normal form gives the stated embedding, and the normal-form arithmetic identifies its sum and product with the surreal operations.
\end{proof}

The proof presupposes the normal-form theorem; it is not a new elementary proof of that theorem. An abstract ordered group $\Gamma$ does not come with a canonical embedding into $\No$ merely because a notation $t^\Gamma$ is written. The canonical embedding above uses the actual inclusion $\Gamma\subseteq\No$.

\subsection{Why the self-Hahn description is not a circular construction}
The slogan
\[
 \No\cong\R((\omega^{\No}))
\]
expresses a representation theorem with \emph{set-sized reverse-well-ordered support}. It does not define $\No$ as the unrestricted set of functions from $\No$ to $\R$. Nor does it mean that a fixed-point equation for an unknown field automatically has a unique solution.

In a class theory, an individual class-Hahn series should be stored by a set support and a set coefficient function on that support. Its zero extension to all surreal exponents is a class function, not a set function with a proper-class domain. In a larger type universe, that zero extension may be an ordinary function type, but the small-support condition still has to be imposed. These are different encodings of the same size discipline.

\subsection{Hahn sums are algebraic operations with admissibility proofs}
A set-indexed family $(a_j)_{j\in J}$ is strongly summable when
\begin{equation}\label{eq:summability}
 \bigcup_{j\in J}\supp(a_j)\text{ is well ordered},\qquad
 \{j:\gamma\in\supp(a_j)\}\text{ is finite for every }\gamma.
\end{equation}
The coefficient at $\gamma$ of $\hSum_j a_j$ is then a finite sum. Both requirements are essential. Well-ordering alone does not allow infinitely many contributions at the same exponent, and local finiteness alone does not prevent a descending support.

The Hahn--Neumann support lemma applies to positive well-ordered supports and validates geometric series, products, and suitable formal substitutions \cite{Neumann}. It is the operative mechanism in the supplied manuscripts, rather than ordinary sequential convergence. In particular, a formal series with arbitrary ordinary coefficients can be evaluated at an infinitesimal even if its ordinary analytic radius is zero. This does not turn it into an ordinary holomorphic germ.

\subsection{A sharp limitation of fixed workspaces}
\begin{proposition}[Universal rescaling fails in a fixed Hahn field]\label{prop:rescaling}
In $K=\C((t^\Q))$, the formal series
\[
 A(X)=\sum_{n\ge0}t^{-n^2}X^n
\]
cannot be evaluated by strong summation at any nonzero $h\in K$. Consequently no nonzero radius in this field has the universal small-scale property asserted for the full class in the analysis manuscript.
\end{proposition}
\begin{proof}
Let $q=v(h)\in\Q$. The leading exponent of $t^{-n^2}h^n$ is $-n^2+nq$. For all sufficiently large $n$, these exponents strictly decrease, because the successive difference is $q-2n-1$. The union of the term supports therefore contains an infinite strictly decreasing sequence and is not well ordered. This violates \eqref{eq:summability}.
\end{proof}

In the full surreal field, one may choose an exponent $\Delta$ dominating the set of coefficient supports and use $\rho=t^\Delta$. The separated-band proof in the supplied analysis manuscript uses precisely this additional freedom \cite[\S4.1]{AnalysisSource}. A formal port to $K_\Gamma$ must either assume an appropriate dominating exponent in $\Gamma$ or enlarge the workspace. It must not assert the full-class theorem uniformly inside every fixed Hahn field.

\subsection{Compatibility under enlargement}
If $\Gamma\hookrightarrow\Delta$ is an order-preserving group embedding, zero extension of coefficients defines the corresponding Hahn-field embedding. It preserves finite operations and those strong sums represented by the same admissible supports. This provides the mathematical basis for coherent workspace enlargement.

Topological compatibility is another question. In $\R((t^\Q))$, $t^n\to0$ in the intrinsic valuation topology. After embedding in the full surreal field, the same sequence does not converge in the full fine topology, whose radii include scales below all its terms. An algebraic and order embedding need not be continuous for these particular source and target topologies. A library should never synthesize the latter property from the former without proof.

\section{Surcomplex numbers require no new large-object axiom}\label{sec:surcomplex}
\subsection{Pairs with complex multiplication}
For an ordered field $F$, define $F[i]$ on pairs by
\begin{align}
 (a,b)+(c,d)&=(a+c,b+d),\\
 (a,b)(c,d)&=(ac-bd,ad+bc),\\
 \overline{(a,b)}&=(a,-b).
\end{align}
This is not the product ring $F\times F$ with coordinatewise multiplication. The latter has zero divisors; the former is the quadratic algebra $F[X]/(X^2+1)$.

\begin{proposition}[Algebra and size of complexification]\label{prop:complexification}
If $F$ is an ordered field, the pair construction above is a field. If $F$ is real closed, it is algebraically closed. Its carrier is a set whenever $F$ is a set, and is a class of set-coded pairs whenever $F$ is a class field of set-coded elements.
\end{proposition}
\begin{proof}
For $(a,b)\ne(0,0)$, ordered-field positivity gives $a^2+b^2>0$, and
\[
 (a,b)^{-1}=\left(\frac{a}{a^2+b^2},\frac{-b}{a^2+b^2}\right).
\]
Direct multiplication gives $(1,0)$. The ring identities are finite polynomial identities. The algebraic-closure statement is the classical characterization of real-closed fields by adjoining a square root of $-1$. Finally, pairing and finite products introduce no new proper-class membership operation.
\end{proof}

There is no ordered-field structure on $F[i]$: every square in an ordered field is nonnegative, whereas $i^2=-1$. Statements about a ``real'' component of a surcomplex number mean an element of $\No$; they must not be confused with ordinary $\R$.

\subsection{The modulus is valued in the base field}
When $F$ is real closed, define
\[
 |a+bi|=\sqrt{a^2+b^2}\in F_{\ge0}.
\]
Multiplicativity follows from the sum-of-two-squares identity. The triangle inequality follows from the corresponding two-dimensional Cauchy--Schwarz identity. These are valid over any real-closed field, at infinite as well as infinitesimal scales. There is no need to reconstruct them specifically from Conway cuts.

For $F=\No$, this modulus is \emph{surreal-valued}, not real-valued. A Lean instance of a real-normed field would not express this modulus. Likewise, replacing an infinite modulus by a real number would change the notion of error tolerance. The geometry and the topology must be formalized using their actual scalar field.

\subsection{Class quotients that do have small individual fibers}
The finite-angle group $\OO_{\R}/(2\pi\Z)$ in the trigonometry manuscript deserves a separate size check. Each coset $\theta+2\pi\Z$ is a set, because it is indexed by $\Z$. The totality of these cosets can therefore be a class of sets. This is unlike the raw-name equivalence classes discussed in \cref{sec:quotients}.

One can avoid the quotient altogether by using the unit-circle group, or choose the unique representative in $[0,2\pi)$. The interval is itself a proper class in the full surreal universe. Boundedness by ordinary reals does not imply sethood: even the interval of infinitesimals contains too many distinct scales to be a set. The chosen interval convention also has an actual branch jump; it should not be mistaken for a globally analytic argument function.

\section{Dependent type theory: small branching and large carriers}\label{sec:types}
\subsection{The universe shift is the central construction, not an inconvenience}
Let $\UU_u$ be a universe of small types. A raw game can have left and right families indexed by arbitrary types $L,R:\UU_u$. A schematic inductive declaration is
\begin{lstlisting}
universe u
inductive RawGame : Type (u + 1) where
  | cut (L R : Type u)
      (left : L -> RawGame) (right : R -> RawGame) : RawGame
\end{lstlisting}
Its carrier is in the next universe because a constructor quantifies over all types in $\UU_u$. The recursive occurrences are strictly positive: they occur as function codomains, not in the domain of a function receiving the entire raw-game type. Well-founded trees of transfinite rank are represented by such inductive or W-type constructions; ``inductive'' does not mean ``finite birthday.''

The associated numeric subtype and its quotient remain in a large universe. They accept small option families but are not themselves a small option family. A type-theoretic version of \cref{thm:nocut} shows that any proposed proof of their smallness would be incompatible with unrestricted small-cut interpolation.

\subsection{An ordinary-inductive route avoids primitive induction-recursion}
The cut presentation appears to define the carrier, its order, and the numeric condition simultaneously. One classical implementation strategy removes this circular appearance:
\begin{enumerate}
\item Define all raw games with small branching by an ordinary inductive construction.
\item Define the game comparison relation by well-founded recursion on pairs of game trees.
\item Define the numeric predicate recursively, requiring numeric options and separation.
\item Prove totality of comparison on numeric games and prove arithmetic congruences.
\item Quotient numeric games by mutual comparison.
\end{enumerate}
The comparison relation on arbitrary games need not be a total order. Totality is a theorem restricted to the numeric subtype. This restriction is indispensable: games and surreal numbers are related objects, not synonymous types.

An alternative is a direct sign-sequence type indexed by a universe of ordinals. That gives a canonical carrier and a straightforward simplicity relation, but the field operations and their properties still require substantial recursion. Neither approach eliminates the mathematical work; it changes where that work is concentrated.

\subsection{Setoids, quotients, and equality}
In intensional Martin-L\"of type theory, an equivalence relation need not be judgmental equality or the identity type. A setoid representation keeps the relation explicitly and requires every operation to respect it. A quotient type packages equivalent representatives into equal outputs, typically with an elimination rule requiring invariance.

For surreal arithmetic, this is not just a syntactic choice. A proof by recursion on a presentation does not automatically define a function on equivalence classes. The invariance theorem must be proved. Conversely, a statement about the canonical birthday of a number must not accidentally use the birthday of an arbitrarily chosen noncanonical representative.

Constructive implementations may prefer setoids to avoid quotient axioms or to expose computation. Classical Lean has native quotients and a large quotient-based algebraic library, so quotienting numeric games is a natural fit. The cost is that reductions involving a representative chosen by choice are generally not computational algorithms for surreal arithmetic.

\subsection{Predicative types, impredicative propositions, and resizing}
Lean's data universes satisfy $\Type_u:\Type_{u+1}$. They are not cumulative: a term of a type in one universe is not automatically a term of the corresponding expression in a larger universe. Explicit lifts such as \texttt{ULift} provide equivalent copies. Universe polymorphism creates a family of declarations parameterized by levels, not an internal type containing every universe \cite{LeanUniverses}.

Lean's $\Prop$ is impredicative, so a proposition may quantify over a large type while remaining a proposition. For example, the statement that every surreal has an additive inverse is a proposition regardless of the universe of its variables. But a proof of a large universal proposition does not shrink the type of its witnesses. A predicate $P:\No_u\to\Prop$ defines a subtype of $\No_u$ at its existing data level; it does not imply that this subtype is $u$-small.

This distinction rules out a tempting error: using the small universe of truth values to encode all surreals as small data. The output of a predicate is a proposition; the domain, graph, subtype, and family of all predicates still have their appropriate universe levels. Resizing must be a stated principle with a suitable scope, not a consequence of impredicativity in $\Prop$.

\section{Constructive set theory and univalent foundations}\label{sec:constructive}
\subsection{IZF, CZF, and the status of order}
IZF retains a power-set style set theory with intuitionistic logic. CZF is designed for predicative constructive mathematics and replaces some classical set-existence principles with constructive collection principles. Type-theoretic interpretations of constructive set theories provide a well-developed connection with dependent types and W-type methods \cite{GambinoAczel}.

These are genuine foundational options, but a classical sign-sequence argument cannot be moved into them by deleting the words ``by excluded middle.'' Classical ordinal comparability, deciding whether an option lies above a point, and turning negated universal statements into witnesses can fail constructively. Even the equivalence
\[
 x<y\quad\Longleftrightarrow\quad\neg(y\le x)
\]
must not be assumed without the relevant order principle. An apartness relation, such as $x\mathrel{\#}y$ meaning $x<y$ or $y<x$, carries positive information absent from mere inequality.

A constructive development should specify its ordinal notion, its inductive or W-type existence principles, its quotient policy, and its use of choice or propositional resizing. We do not claim that bare CZF proves every desired large inductive construction, nor that all constructive surreal variants are equivalent to the classical sign-sequence class without additional logical hypotheses.

\subsection{The higher inductive-inductive construction}
The HoTT book gives a higher inductive-inductive presentation of surreals \cite[\S11.6]{HoTT}. One simultaneously defines the type $\No$, its non-strict order, and its strict order. The carrier has a cut constructor for small separated left and right families. It also has a path constructor identifying $x$ and $y$ when $x\le y$ and $y\le x$.

Schematically, the order rules read
\begin{align*}
 [\forall x^L\;(x^L<y)]\ \land\ [\forall y^R\;(x<y^R)]&\Longrightarrow x\le y,\\
 [\exists x^R\;(x^R\le y)]\ \lor\ [\exists y^L\;(x\le y^L)]&\Longrightarrow x<y.
\end{align*}
Here the notation refers to specified cuts presenting $x$ and $y$. The full construction must supply the simultaneous induction rules and account for the identity and propositionality conditions. It is not merely an ordinary inductive datatype with a suggestive equation attached.

Shulman's discussion explains why the constructive definition treats strict and non-strict comparison separately and relates it to plump ordinals \cite{Shulman}. Joseph develops an ordered commutative-ring structure with apartness in univalent foundations \cite{Joseph}. These sources establish that the constructive approach is mathematically substantial. They do not, by themselves, constitute a Lean 4 implementation of the complete surcomplex analytic theory.

\subsection{Why one cannot simply add ordinary univalence to Lean}
Lean's equality is a proposition and its proofs are proof-irrelevant. Consequently, for fixed types $A,B$, the identity type $A=B$ has at most one proof. In a univalent universe, identity between types corresponds to equivalence between them, retaining the nontrivial automorphisms of a type.

\begin{proposition}[A basic incompatibility]\label{prop:univalence}
A universe with proof-irrelevant equality of types cannot satisfy the usual univalence equivalence $(A=B)\simeq(A\simeq B)$ for all its types if it contains a two-element type with its two distinct automorphisms.
\end{proposition}
\begin{proof}
Take $A=B$ to be the two-element type. Proof irrelevance makes $A=A$ a subsingleton. An equivalent type would also be a subsingleton. But $A\simeq A$ contains the identity automorphism and the automorphism exchanging the two elements; evaluating them at either element proves they are different.
\end{proof}

This does not preclude studying HoTT \emph{inside} Lean by formalizing its syntax or a model. It does preclude identifying HoTT's univalent universes with Lean's ordinary types and adding the usual univalence statement unchanged. Propositional extensionality is a different, much more restricted assertion and is compatible with Lean's foundation \cite{LeanAxioms}.

For the requested classical manuscripts, native Lean quotients and universe-indexed games are therefore a more direct route than importing a higher inductive-inductive axiom package. A constructive or univalent project is worthwhile for different goals, but it must explicitly re-establish the field, order, and analytic hypotheses it needs.

\subsection{Categorical and algebraic-set-theoretic perspectives}
A category equipped with a notion of small maps can distinguish small fibers from a large ambient object; this provides a conceptual analogue of small option families and small supports. Algebraic approaches to surreal structures have been developed in a class-set-theoretic setting \cite{RangelMariano}. Such frameworks can organize universal constructions and comparison functors elegantly.

They do not remove the diagonal obstruction. Internally in any logic where $<$ is irreflexive, a rule producing a strict upper bound for every subobject of the same carrier fails on its maximal subobject. A categorical account must restrict the small maps or the admissible cuts. Its semantics, quotient existence, and inductive constructions should be stated explicitly rather than hidden behind the word ``universal.''

\section{Axiomatizing the structure, not only constructing it}\label{sec:axioms}
\subsection{Real closedness is an interface, not a characterization of \texorpdfstring{$\No$}{No}}
The axioms of real-closed ordered fields capture exactly the algebra needed for many finite geometric arguments. They do not characterize the surreal class: $\R$ and numerous non-Archimedean fields satisfy them as well. Adding an infinite element only rules out Archimedean examples. Adding a specified family of infinitesimals still does not recover Conway's simplicity structure.

The same warning applies to an abstract algebraically closed field advertised as ``surcomplex.'' It may suffice for polynomial factorization, but it does not supply a distinguished conjugation, an ordered fixed field, a Conway monomial map, or an ordinary complex subfield. Those are additional data with laws.

\subsection{A safe abstract package}
A useful classical abstract interface separates the following ingredients:
\begin{enumerate}
\item An ordered-field carrier together with an ordinal-valued birthday and a simplicity relation having set-sized predecessor chains.
\item A full sign-tree representation: each element has a unique ordinal-length sign sequence, every such small or set-sized sequence is represented, and numerical order is the termination-sensitive lexicographic order.
\item A cut operation for separated small families, with its simplest-element property and a rank bound.
\item Arithmetic compatible with the recursive Conway rules, not just arbitrary field operations on the same carrier.
\end{enumerate}
This is a structural specification, not a proposed minimal finite axiom system. In a class theory the quantification over sign sequences is quantification over set codes; in a universe-indexed theory it is restricted to the stated universe.

\begin{proposition}[Rigidity of the full sign-tree specification]\label{prop:rigidity}
Two realizations of the full sign-tree specification over the same ordinal and set universe have a unique isomorphism preserving sign sequences. If both use the recursive Conway arithmetic, that isomorphism preserves the arithmetic.
\end{proposition}
\begin{proof}
Map each element to the unique element in the second realization with the same sign sequence. Existence and uniqueness of represented sequences make this a bijection; the order and simplicity relations are preserved by their definitions. Its compatibility with addition and multiplication follows by well-founded induction on the input presentations, since it identifies corresponding option values and corresponding simplest cuts. No other sign-preserving map can differ on any element.
\end{proof}

Ehrlich's work supplies more conceptual simplicity-hierarchical axiomatizations and initial-embedding results \cite{Ehrlich,EKHierarchy}. Their strength comes from the extra tree structure and the specified notion of initial embedding, not merely from field axioms plus a vague maximality condition. In particular, maximality in a structured category must not be confused with the absence of every possible ordered-field extension.

\subsection{A current development, with its status kept explicit}
A March 2026 announcement by Hamkins describes joint work with Junhong Chen and Ruizhi Yang on a first-order theory of surreal arithmetic using the surreal field equipped with its birthday-order structure. The announcement states a bi-interpretability result with ZFC and explicitly describes the project as work in progress \cite{SurrealArithmetic}. We report that announcement as such; we do not use it as a proved dependency of this article or as a complete published axiomatization supplied here.

The relevant point is the role of the enriched structure. One should not read the announcement as a claim about the pure language of real-closed ordered fields. Birthday and simplicity information can encode far more than polynomial field identities.

\subsection{Analytic canonicity depends on the chosen category}\label{subsec:canonicalexp}
The trigonometry manuscript correctly separates finite-angle trigonometry from extensions determined only by weak global group and local differential axioms. Outside literature adds an important qualification: Ehrlich and Kaplan construct canonical surcomplex exponentials using an \emph{initial integer part} \cite[\S11]{EKExp}. This is extra structural information, not a contradiction of nonuniqueness under weaker hypotheses.

In the full surreal field, the omnific integers have the normal-form description
\[
 \Oz=\Pi\oplus\Z,
\]
where $\Pi$ consists of purely infinite normal forms. For finite-angle $\cis$, put
\[
 U_0(p+u)=\cis(u),\qquad p\in\Pi,\quad u\in\OO_{\R}.
\]
Its kernel is $\Pi\oplus2\pi\Z=2\pi\Oz$, because multiplication by the nonzero ordinary real $2\pi$ preserves $\Pi$. Thus
\[
 E_0(x+iy)=\exp_{\No}(x)U_0(y)
\]
has kernel $2\pi i\Oz$. This is the normal-form description consistent with the canonical initial-integer-part extension in the cited theorem. Specifying that integer part kills the free purely infinite phase component.

\begin{proposition}[Why an infinite period is not a paradox]\label{prop:period}
Let $U:(\No,+)\to\mathbb T(\No)$ be a homomorphism extending the finite-angle map, and assume that the latter is already onto the full unit circle. Then every coset of the finite additive subgroup contains an element of $\ker U$.
\end{proposition}
\begin{proof}
Given $x\in\No$, choose a finite $a$ such that $\cis(a)=U(x)^{-1}$. Then $U(x+a)=1$. If $x$ is outside the finite subgroup, $x+a$ is still infinite.
\end{proof}

Accordingly, the kernel $2\pi i\Oz$ is compatible with, and illustrates, the manuscript's unavoidable-infinite-period theorem. A formal library should expose both levels: finite-angle identities requiring no infinite-phase choice, and the extra integer-part or character data used by a selected global extension. Neither the foundation nor the bare field axioms determine that data for free.

\section{A size-correct formalization of the fine topology}\label{sec:topology}
\subsection{The small positive lower-bound principle}
The fundamental property used by the analysis manuscript is
\begin{equation}\label{eq:lowerbound}
 A\subseteq\No_{>0}\text{ a set}\quad\Longrightarrow\quad
 \exists\delta>0\ \forall a\in A\;(\delta<a).
\end{equation}
For nonempty $A$, the cut $\{0\mid A\}$ proves this; the empty case is immediate. Notice that $\delta$ may not belong to any previously specified subfield containing $A$.

In a universe-indexed construction, the corresponding statement has the essential hypothesis that $A$ is small at the option-family level. It does not hold for the entire higher-universe type of positive surreals. That type contains a positive element smaller than any proposed positive lower bound.

\begin{theorem}[Small subsets and small-index nets]\label{thm:smalltopology}
Assume \eqref{eq:lowerbound} for the permitted small families. Give $F[i]$ the ball topology with radii in $F_{>0}$, where $F$ is the corresponding surreal field or universe-relative carrier. Then every permitted small subset is closed and discrete. Every convergent net with permitted small index type is eventually equal to its limit. Every Cauchy net with permitted small index type is eventually constant.
\end{theorem}
\begin{proof}
Let $A$ be a small subset and fix $p$. The nonzero distances $|p-a|$ for $a\in A$ form a small family. Choose $\delta>0$ smaller than all of them. The ball $B(p,\delta)$ meets $A$ in at most $\{p\}$. This isolates each member and separates each exterior point from $A$, proving discreteness and closedness.

For a net $z_j\to z$, apply the same argument to its nonzero distances from $z$. Eventual membership in the resulting ball forces $z_j=z$. For a Cauchy net use the small family of nonzero pairwise distances $|z_j-z_k|$. A Cauchy tail has all these distances less than $\delta$, so its values are equal.
\end{proof}

This is the precise mathematical content of the manuscript's set-sized topological degeneracy theorem \cite[\S3]{AnalysisSource}. A Lean theorem quantifying over arbitrary \texttt{Set Surreal} without a smallness hypothesis would say something stronger and false.

\subsection{An explicit larger-index counterexample}
Let $D=F_{>0}$ and direct it by reverse numerical order: $r\preceq s$ means $s\le r$. Define $x_r=r$. For any $\epsilon>0$, once $r\preceq s$ with $r=\epsilon/2$, one has $0<x_s=s\le\epsilon/2<\epsilon$. Hence this net converges to zero and is never zero.

In a class theory, $D$ is a proper class and the net is a class-indexed net. In Lean, for $F=\No_u:\Type_{u+1}$, $D$ is a perfectly legitimate type in the larger universe, but it is not a permitted $u$-small index type. There is no conflict with \cref{thm:smalltopology}. In particular, replacing ``sequence'' by ``net'' does not fix a missing size annotation unless the index universe is also specified.

\subsection{Why ordinary real-normed analysis is not a drop-in replacement}
The full fine topology is not discrete: every ball around zero contains half its radius. Yet its small subsets are discrete. This distinction prevents importing the ordinary intuition that pointwise isolation on all small samples makes the whole space discrete.

It also prevents a real metric from representing this topology in the usual small-index setting. A countable local basis would have a positive radius inside each of its members. Applying \eqref{eq:lowerbound} to those radii produces a smaller ball that contains none of those basis neighborhoods: each chosen radius can be halved to produce a point still in its own ball but outside a sufficiently smaller common ball. More explicitly, first choose $\delta<r_n/2$ for every $n$; then $r_n/2\notin B(0,\delta)$. Thus no countable ball basis is coinitial. Ordinary real metrics always have a countable local basis.

A correct implementation may use \texttt{TopologicalSpace} and filters, because these are not restricted to real metrics. Function differentiation can be defined by
\begin{equation}\label{eq:derivative}
 \forall\epsilon\in F_{>0}\ \exists\delta\in F_{>0}\ \forall h,
 \quad 0<|h|<\delta\Rightarrow
 \left|\frac{f(z+h)-f(z)}h-d\right|<\epsilon.
\end{equation}
Here $f$ is defined on a neighborhood of $z$, and the chosen ball of increments is required to stay in its domain. The derivative is unique by ordered-field estimates. This definition does not require a real-valued norm, a sequence limit, or a derivation on the scalar field.

\subsection{Fine-local functions and coherent functions remain different}
A function that is constant on each infinitesimal coset can be nonconstant globally and have derivative zero everywhere. Thus the usual implication ``zero derivative on a connected domain implies constant'' cannot simply be reused with an ordinary connected base domain and a class-sized surreal thickening. The supplied analysis manuscript imposes common-domain Hahn coherence precisely to obtain stronger global conclusions.

In a formalization, \texttt{FineHasDerivAt}, a formal-series germ, a Hahn-coherent datum, and a scalar-field derivation should therefore be distinct types or predicates. Proving implications among them is useful; identifying them by notation is not.

\section{Lean: the current implementation landscape}\label{sec:leanstatus}
\subsection{A pinned source inspection, not an execution claim}
The public \texttt{vihdzp/combinatorial-games} repository was inspected at the following revision on September 21, 2026 \cite{LeanGames}:
\begin{center}\small
\begin{tabular}{@{}L{0.22\textwidth}L{0.71\textwidth}@{}}
\toprule
Item & Recorded value\\\midrule
Repository revision & \texttt{02b4a908ea2ecfefffecb438f691951a814a5264}\\
Commit date & September 19, 2026\\
Lean toolchain & \texttt{leanprover/lean4:v4.35.0-rc2}\\
Mathlib revision & \texttt{dec5b2b780537b6eaf7f5e5f000c12f7387fb24d}\\
\bottomrule
\end{tabular}
\end{center}
The toolchain and mathlib revisions were read from \emph{that pinned commit}, not inferred from search-engine snippets or from an unrelated cached branch page. No Lean compiler was available in the execution environment used to prepare this article. Source inspection, the mathematical proofs here, and the illustrative files must therefore not be described as a completed local build or an axiom-checked verification of the manuscripts.

\subsection{What the inspected code establishes as an implementation fact}
The basic surreal carrier is defined by antisymmetrizing numeric games:
\begin{lstlisting}
def Surreal : Type (u + 1) :=
  Antisymmetrization (Subtype Numeric) (· ≤ ·)
\end{lstlisting}
The same module supplies the quotient map, numerical equality of representatives, a linear order, additive structure, and a choice-based representative operation. Its set-based cut interface uses smallness conditions on the option families. This is the universe-sensitive pattern developed above, not a single small type accepting arbitrary subsets of itself.

The division module contains an actual \texttt{Field Surreal} instance. Its proof constructs inverse numeric games and proves inverse correctness by well-founded induction. Thus historical documentation reporting that surreal multiplication or field structure has not been implemented should not be used as a description of this inspected revision. Conversely, a prose phrase such as ``complete ordered field'' in a source comment must not be interpreted as a Dedekind-completeness instance; \cref{sec:obstructions} explains why that mathematical claim would be wrong.

The inspected \texttt{Surreal/HahnSeries/Basic.lean} defines \texttt{SurrealHahnSeries} as a field of small-support Hahn series over the order-dual surreal exponents. Its constructor takes a coefficient function, a proof of small support, and reverse well-foundedness. This is directly relevant to the manuscripts' support discipline. It does not follow merely from that type's existence that an ordered-field equivalence with the game-based carrier, a full normal-form theorem, real closedness, or Gonshor exponentiation has been established by the modules inspected here. Those remain \emph{unverified dependencies in this audit}, not assertions of absence from every available project.

\subsection{Mathlib already distinguishes formal and topological sums}
Mathlib's Hahn-series API has \texttt{HahnSeries.SummableFamily} and its formal sum \texttt{hsum}. Its documentation explicitly distinguishes this notion from topological summability \cite{MathlibHahn}. This is the correct starting point for \eqref{eq:summability}; inventing a new use of ordinary \texttt{tsum} would obscure an already-supported distinction.

The general Hahn API permits ordered structures more general than a linearly ordered group and uses partially well-ordered support in its general formulation. A surreal specialization should prove that these general hypotheses yield the required well-ordered supports in the chosen linear order. It should also establish smallness of the resulting support, not only its order-theoretic shape.

Mathlib also has an \texttt{IsRealClosed} interface \cite{MathlibRealClosed}. An interface name is not an instance for a particular type. Before depending on real closedness of the actual game-based \texttt{Surreal}, the project must locate or prove the relevant instance and check its dependencies. Similar care is needed when deriving algebraic closedness of its quadratic extension.

\subsection{An internal model of set theory is an alternative, not a necessity}
Mathlib provides \texttt{ZFSet : Type (u + 1)}, constructed from pre-sets modulo extensional equivalence, together with a class theory and set operations \cite{MathlibZFC}. This permits a formalization \emph{inside an internal set-theoretic model}: define surreal sign sequences as internal sets, prove internal class statements, and compare the result with the native game-based type.

This approach is well suited to proving interpretation or conservativity results. It is often less convenient for ordinary analysis, because each ordinary field, polynomial, function, and set operation may need an internal-to-native translation. The existence of a ZFC model in Lean also does not mean that Lean is simply ZFC under another spelling; the construction uses Lean's ambient type theory and universe hierarchy.

A terminological warning is especially important here: the library's technical predicate \texttt{ZFSet.Definable}, and its theorem \texttt{Classical.allZFSetDefinable}, should not be read as saying that every external function is definable by a first-order formula in the language of set theory. The API concerns representative-level realizability of functions on the quotient. First-order definability for a metamathematical interpretation must be formalized with actual syntax and satisfaction.

\section{A proposed Lean architecture for the two manuscripts}\label{sec:architecture}
\subsection{Layer A: reusable finite algebra}
Start with an arbitrary ordered field $F$, not a concrete surreal datatype. Define its quadratic extension with the correct cross-term multiplication. Prove conjugation, norm-square identities, inverses, and the Gram identity. Add real-closedness only where positive square roots or polynomial factorization are genuinely required.

This layer handles much of the geometric algebra in the trigonometry manuscript: determinant identities, the cosine-law polynomial identity, Heron's factorization, circumcenter equations, and Ptolemy's inequality once an appropriate modulus is available. A rational parametrization of the unit circle also needs only field algebra and order:
\[
 t\longmapsto\frac{1-t^2}{1+t^2}+i\frac{2t}{1+t^2}.
\]
Representing an abstract direction by a circle point does not yet assert that a chosen sine function parametrizes it by numerical finite angles. That is a later theorem.

\subsection{Layer B: set-sized Hahn analysis}
Use $F_\Gamma$ and $K_\Gamma$ as ordinary types. Make the ordered value group, its divisibility when needed, and the coefficient field explicit parameters. Reuse the existing formal-summability API. Establish standard part for finite series from the coefficient at exponent zero, and distinguish finite elements from positive-valuation infinitesimals.

Implement evaluation as an operation whose input includes an admissibility proof, or as a total operation with an accompanying theorem guaranteeing its mathematical meaning on admissible inputs. The former style is particularly clear:
\[
 \operatorname{evaluate}(A,h,\text{proof that }(a_nh^n)_n\text{ is strongly summable}).
\]
An ordinary polynomial needs no infinite-sum proof. A formal power series does. Keeping these constructors distinct prevents accidental evaluation outside the proven domain.

The next targets are the positive-support geometric series, binomial series, infinitesimal exponential and logarithm, and Taylor lifting of ordinary analytic germs. At a finite point $r+\epsilon$, first evaluate the ordinary function at $r$, then substitute $\epsilon$ in its Taylor series. Do not replace this by the ordinary Maclaurin series at the nonzero real $r$: infinitely many of those terms can occupy the same Hahn exponent.

\subsection{Layer C: support-controlled coherent functions}
For an ordinary open $U\subseteq\C$, represent a coherent datum as a Hahn series with coefficients in the ring $\OO(U)$ of ordinary holomorphic sections. Equality of sections is equality on $U$; irrelevant values outside $U$ should not become part of the mathematical object.

Differentiation of coefficients does not enlarge the common support. Integration along a fixed ordinary admissible contour also acts coefficientwise and does not enlarge it. These observations explain why ordinary complex-analytic theorems can be applied to the coefficients and then reassembled. They do not prove arbitrary composition: for that, one must control the supports of the inner perturbations and verify that the ordinary part maps into the common coefficient domain.

A useful composition interface has an ordinary base map $g_0:V\to U$ and an infinitesimal perturbation whose coefficient support lies in a common positive well-ordered set $T$. Taylor expansion then produces supports in a set of the form $S+T^*$, where $S$ is the outer support. The formalization must prove local finiteness of coefficient contributions as well as well-ordering. The same positive-support recursion is the natural implementation route for the cluster-factor lifting in the trigonometry manuscript.

\subsection{Layer D: the universe-indexed surreal bridge}
The bridge should establish explicit maps among numeric games, canonical sign sequences, and small-support normal forms. Its main deliverables are preservation of order, field operations, ordinary real coefficients, Conway monomials, and admissible sums. The strongest comparison theorem should be stated with its support universe visible.

A generic \texttt{HahnSeries Surreal R} type, without a small-support restriction, can contain series too large to represent a surreal at the chosen level. Consider the family of singleton monomials indexed by every ordinal of the smaller universe. Each individual support is small and contributions are disjoint, but their union is not small at that level. A class version would be the inadmissible sum
\[
 \sum_{\alpha\in\On}\omega^{-\alpha}.
\]
A proof that the union is well ordered does not prove its sethood. This is exactly the reason a \texttt{SurrealHahnSeries}-style small-support subtype is needed.

\subsection{Layer E: topology, functions, and enrichment}
Define the fine topology and derivative using all radii of the chosen carrier. State the discreteness and eventual-constancy theorems only for small subsets or index types. Workspace embeddings get separate algebraic, order, summability, and topological compatibility theorems; the last one is not automatic.

After the normal-form bridge, add canonical finite sine and cosine, the finite-angle quotient, and the circle logarithm. Add Gonshor's real exponential as a separately proved structure. The strip exponential then follows from finite-angle trigonometry and the real exponential. For a global exponential, require a specified character or the stronger initial-integer-part structure discussed in \cref{subsec:canonicalexp}.

Do not package all of this into a single enormous typeclass at the beginning. A theorem about Heron's identity should not acquire an accidental dependency on a global exponential or class-choice axiom. Small interfaces make both reuse and foundational auditing substantially easier.

\subsection{Suggested module boundaries and completion criteria}
\begin{center}\small
\begin{tabular}{@{}L{0.27\textwidth}L{0.64\textwidth}@{}}
\toprule
Proposed module group & A meaningful completion test\\\midrule
\texttt{Algebra/Quadratic} & Pair multiplication, conjugation, inverse, norm-square identities.\\
\texttt{Hahn/Support} & Strong sum API, positive-support substitutions, smallness preservation.\\
\texttt{Hahn/Germs} & Formal evaluation and recentering with explicit domain witnesses.\\
\texttt{Hahn/Coherent} & Coefficientwise identity, differentiation, and contour theorems.\\
\texttt{Surreal/Bridge} & Normal-form comparison with algebra and order compatibility.\\
\texttt{Surreal/FineTopology} & Small-index statements plus the larger-index counterexample.\\
\texttt{Surcomplex/Angles} & Full unit-circle parametrization and finite-angle kernel.\\
\texttt{Surcomplex/Enrichment} & Real exponential, strip functions, specified global phases.\\
\bottomrule
\end{tabular}
\end{center}
These names are an architectural proposal, not claims that modules with these exact paths already exist. A dependency gate that has not yet been proved should remain a clearly stated hypothesis in downstream work, not be replaced with a new unexplained axiom called a theorem.

\section{Translating the manuscripts: a theorem-by-theorem dependency audit}\label{sec:audit}
The following audit concerns foundational and representation requirements. It does not certify all analytic proof details in the original texts.

\begin{longtable}{@{}L{0.25\textwidth}L{0.32\textwidth}L{0.34\textwidth}@{}}
\toprule
Source result family & Sufficient working layer & Obligation that must remain explicit\\\midrule
\endfirsthead
\toprule
Source result family & Sufficient working layer & Obligation that must remain explicit\\\midrule
\endhead
\bottomrule\endfoot
Finite vector and triangle identities & Generic ordered or real-closed field; quadratic extension & Positive square roots and nonzero denominators; genuine angle interpretation added separately.\\[0.4em]
Strong sums and infinitesimal Taylor lifts & A set-sized Hahn workspace & Well-ordered union, finite coefficient incidence, and correct ordinary base value.\\[0.4em]
Polar decomposition by finite angles & Hahn infinitesimal exponential/logarithm and ordinary circle & Surjectivity on every unit point, exact kernel, and quotient size.\\[0.4em]
Arbitrary-coefficient small-radius evaluation & Full class, or universe-relative small-family boundedness & A dominating exponent may lie outside the initial workspace; \cref{prop:rescaling}.\\[0.4em]
Fine-topology degeneracy & Small positive lower-bound principle & Smallness of the subset or net index, not only the existence of a Lean type.\\[0.4em]
Coherent Cauchy and identity theorems & Ordinary holomorphic coefficients on one common domain & Support compatibility, ordinary domain hypotheses, and coefficientwise contour definition.\\[0.4em]
Preparation and root-cluster conservation & Finite polynomial/algebraic data plus support recursion & Algebraic closure; finite coefficient equations; root supports may enlarge beyond factor supports.\\[0.4em]
Finite Fourier and Fej\'er--Riesz results & Finite sums and real-closed-field algebra, with circle interpretation & Ordinary finite degree; positivity on the entire relevant circle, not only its real shadow.\\[0.4em]
Strip exponential and hyperbolic geometry & Gonshor exponential plus finite phases & The real exponential structure is additional to the bare field interface.\\[0.4em]
Global phase classification & Additive normal-form splitting and a specified class/type homomorphism & Quantify over characters at the right sort; do not form an illegal class of classes.\\[0.4em]
All-scale rigidity & A full class or universe-relative quantification over scales & No inference from a single fixed workspace; each coefficient family must be small at the operative level.\\[0.4em]
\end{longtable}

\subsection{A sound localization argument has two halves}
To reduce a theorem to a workspace, first prove that all input data fit in it, and then prove that every construction used by the proof stays there or specify a controlled enlargement. \Cref{thm:workspace} supplies the first half, not automatically the second. Polynomial factorization stays in a divisible algebraically closed Hahn workspace. A rescaling that dominates every exponent in a coefficient sequence need not.

Similarly, the image of a coherent datum under differentiation or an ordinary contour functional has support contained in the original support. A statement quantifying over \emph{all} surreal scales is not thereby reduced to one such datum. The universal quantifier has to be preserved in the formal statement.

\subsection{Finite operations and infinite conclusions}
Much of the useful geometry is finite algebra even when lengths are infinite. The computational complexity of evaluating a particular surreal coefficient does not affect the logical validity of a polynomial identity. This makes generic algebra an unusually high-return first formalization target.

By contrast, validating finitely many Taylor coefficients cannot verify a support theorem, a transfinite recursion, or an assertion about every scale. The symbolic checks reported in the trigonometry manuscript are useful sanity checks of specific identities, not substitutes for these proofs \cite[\S17.3]{TrigSource}. The same limitation applies to the illustrative files in this archive: they are intentionally much narrower than the complete analysis project.

\section{Relative consistency and proof-assistant trust}\label{sec:trust}
\subsection{What a construction can guarantee}
If a base theory $T$ proves that a definable structure satisfies a collection of axioms $A$, then the interpretation is a relative consistency justification: a contradiction obtained from the interpreted mathematics would yield a contradiction in $T$. Definitional extensions do not add new mathematical existence assumptions beyond what their definitions and existence proofs establish.

This is not an absolute proof of $\Con(T)$. In particular, a strong effectively axiomatized theory cannot generally prove its own consistency under the usual hypotheses of the second incompleteness theorem. A Lean proof of a model or an interpretation also has the foundation and universe assumptions of its ambient system. These qualifications are not reasons to distrust surreal arithmetic; they are the same foundational qualifications that apply to ordinary formalized mathematics.

\subsection{Kernel checking and mathematical faithfulness are separate}
Lean's kernel checks proof terms against definitions, inference rules, and admitted axioms. Tactics that synthesize such terms are not substitutes for the kernel. A final theorem can nevertheless be mathematically misidentified if its formal statement uses the wrong notion of summability, omits a smallness hypothesis by strengthening the base assumptions, or proves a result only for a restricted surrogate field.

A reliable verification report therefore records both a successful build and a statement-level correspondence with the intended theorem. The report should say whether the carrier is $\No_u$, a birthday fragment, a Hahn workspace, or an abstract real-closed field, and whether a function is fine-local, Hahn-coherent, or equipped with a chosen global phase.

\subsection{Axioms should be audited, not hidden}
Lean provides \texttt{\#print axioms} for declaration dependencies. The axiom \texttt{sorryAx} can prove arbitrary statements and does not belong in a finished proof \cite{LeanAxioms}. Classical choice and propositional extensionality are standard explicit ingredients of Lean's classical mathematical environment, not silent surreal-specific assumptions. Quotient principles and native-computation trust mechanisms should likewise be understood in the environment where the project is built.

An axiom listing is not a complete description of the kernel's rules or the semantic strength of its universe hierarchy. Nor does an empty list establish the consistency of the kernel. It is, however, a valuable way to detect unproved placeholders or unexpected project-specific axioms in a purported theorem.

The supplied archive contains no claim of an executed Lean build. Its example files are labeled as uncompiled illustrations, and the source audit records exactly which repository declarations were inspected. The PDF itself was compiled from the accompanying LaTeX source.

\section{Comparing the approaches and choosing a foundation}\label{sec:comparison}
\begin{longtable}{@{}L{0.20\textwidth}L{0.34\textwidth}L{0.37\textwidth}@{}}
\toprule
Approach & Main advantage & Cost or boundary\\\midrule
\endfirsthead
\toprule
Approach & Main advantage & Cost or boundary\\\midrule
\endhead
\bottomrule\endfoot
ZF/ZFC definable classes & Classical carrier and operations without making $\No$ a set; no inaccessible needed merely for the construction. & Class quantification becomes a definability schema; every recursive graph and set-bounding step must be justified.\\[0.4em]
GB/GBC/NBG & Natural language for particular global functions and full-class statements; closest to the manuscripts. & Classes cannot be elements; elementary comprehension and class recursion have limits; choice convention must be specified.\\[0.4em]
KM or stronger class theories & Supports stronger class comprehension and recursion. & Stronger foundational commitment; still no unrestricted cuts or naive hyperclasses.\\[0.4em]
Birthday fragments & Concrete set-sized construction stages and clear ordinal bookkeeping. & An arbitrary stage need not be a field; analytic closure requires specific theorems.\\[0.4em]
Set-sized Hahn fields & Direct, modular access to support-controlled algebra and analysis. & Not the entire surreal universe; full-class rescaling and topology do not automatically survive restriction.\\[0.4em]
Grothendieck universes / TG & Clean internal-small versus external-large distinction and rich set operations. & Set-theoretic universe existence adds strength; relative totality is not absolute totality.\\[0.4em]
Classical dependent types & Native arithmetic libraries, quotients, well-founded recursion, explicit universe polymorphism. & Universe transport and smallness proofs are mathematical obligations; quotient representatives are not computational normal forms.\\[0.4em]
Constructive set/type theory & Exposes witnesses, apartness, and computational content. & Classical order arguments and some existence principles must be replaced or added explicitly.\\[0.4em]
HoTT / higher induction & Numerical equivalence can be built into identity; constructive order is handled simultaneously. & Needs an appropriate higher inductive-inductive foundation; not a plug-in axiom for Lean's ordinary equality.\\[0.4em]
Internal set theory in Lean & Suitable for formal comparison of foundations and internal models. & Additional encodings and bridges to native algebra and analysis; not needed for every application.\\[0.4em]
\end{longtable}

For the requested manuscripts, the recommended mathematical foundation is a precisely specified class theory or an explicitly definable-class ZFC presentation. The recommended \emph{implementation} foundation is a universe-indexed classical Lean development with a set-sized Hahn analytic core. These recommendations are compatible: one is a convenient mathematical language, while the other is a proof-engineering representation with a documented interpretation.

The unifying rule is simple but has substantial consequences: \emph{an individual object can have small data without the totality of all such objects being small}. Applied consistently, it explains why surreal cuts, surcomplex fields, normal forms, and strong sums are legitimate, why full-class analysis differs from analysis in one field, and why a formalization must preserve the size annotations rather than erase them.

\appendix
\section{Illustrative Lean files and reproducibility notes}\label{app:lean}
\subsection{A logical guard against the unrestricted cut error}
The accompanying \texttt{lean/LogicalGuards.lean} contains this small proof pattern:
\begin{lstlisting}
import Mathlib
universe u

theorem noUnrestrictedUpperBounds (F : Type u) [Preorder F] :
    ¬ (∀ s : Set F, ∃ x : F, ∀ y ∈ s, y < x) := by
  intro h
  obtain ⟨x, hx⟩ := h Set.univ
  exact (lt_irrefl x) (hx x (Set.mem_univ x))
\end{lstlisting}
It is a useful negative test for a proposed abstract interface. The file also includes the raw-game universe skeleton and polynomial identities for quadratic pairs. These are illustrative source files, not a compiled implementation of surreal numbers.

\Needspace{16\baselineskip}
\subsection{A source-based API probe}
The accompanying \texttt{lean/SourceAudit.lean} imports the inspected modules and lists basic probes:
\begin{lstlisting}
import CombinatorialGames.Surreal.Division
import CombinatorialGames.Surreal.HahnSeries.Basic

set_option pp.universes true
#check Surreal
#check SurrealHahnSeries
#check SurrealHahnSeries.mk
#synth Field Surreal.{0}
#synth Field SurrealHahnSeries.{0}
#check Surreal.mk_eq_mk
#print axioms Surreal.mk_eq_mk
\end{lstlisting}
The probe does not assert that real closedness, a normal-form equivalence, or analytic extensions can be synthesized. Those are deliberately separate audit questions. A successful run would confirm only the displayed probes under the selected dependency configuration.

To reproduce the intended environment, use the repository at the recorded commit, retain its \texttt{lean-toolchain} and \texttt{lake-manifest.json}, and run the files with that project's \texttt{lake env lean}. Do not run a dependency update and then describe the result as a test of the pinned revision. The README supplies the exact commands and labels their execution status.

\section{A compact list of forbidden shortcuts}\label{app:shortcuts}
\begin{enumerate}
\item Do not apply the cut operation to all surreals at its own operative size level.
\item Do not treat a proper-class equivalence class as an element of an NBG class.
\item Do not identify surreal cut interpolation with Dedekind completeness.
\item Do not use an arbitrary birthday bound as a field instance without closure proofs.
\item Do not identify the self-Hahn normal-form theorem with an unrestricted function-space definition.
\item Do not confuse well-ordered support with small support, or either with finite coefficient incidence.
\item Do not replace strong Hahn summation by a topological \texttt{tsum}.
\item Do not infer workspace stability from the fact that the inputs lie in a workspace.
\item Do not remove the small-index qualification from the fine-net theorem.
\item Do not substitute a real-valued norm for the actual surreal-valued modulus.
\item Do not derive global analytic uniqueness from a local differential equation alone.
\item Do not call illustrative source, symbolic checks, or a module comment a machine-checked proof of the intended theorem.
\end{enumerate}
These restrictions are not ad hoc prohibitions. Each records a specific size, domain, algebraic, or verification obligation explained in the article.

\clearpage
\begingroup\small\setlength{\parskip}{0pt}
\begin{thebibliography}{99}
\setlength{\itemsep}{0.45em}\raggedright
\addcontentsline{toc}{section}{References}
\bibitem{AnalysisSource}
\emph{Surcomplex Analysis: Infinitesimal Calculus, Hahn-Coherent Holomorphy, and Contour Theory}.
Supplied manuscript, \texttt{surcomplex\_analysis(2).tex}, September 21, 2026. Foundational conventions in \S2; fine topology in \S3; small-scale evaluation in \S4; coherent functions and global-scale theory in the later sections. Not independently machine-verified here.
\bibitem{TrigSource}
\emph{Trigonometry on the Surcomplex Plane: Canonical Angles, Infinite Triangles, Infinitesimal Contact, and Hahn-Analytic Oscillation}.
Supplied manuscript, \texttt{surcomplex\_trigonometry.tex}, September 21, 2026. Size and support conventions in \S2, angle group in \S4, global phases in \S16. Not independently machine-verified here.
\bibitem{Conway}
John H. Conway. \emph{On Numbers and Games}. Academic Press, 1976; second edition, A K Peters, 2001.
\bibitem{Gonshor}
Harry Gonshor. \emph{An Introduction to the Theory of Surreal Numbers}. London Mathematical Society Lecture Note Series 110, Cambridge University Press, 1986.
\href{https://doi.org/10.1017/CBO9780511629143}{doi:10.1017/CBO9780511629143}.
\bibitem{Alling}
Norman L. Alling. \emph{Foundations of Analysis over Surreal Number Fields}. North-Holland Mathematics Studies 141, North-Holland, 1987.
\bibitem{vdDE}
Lou van den Dries and Philip Ehrlich. Fields of surreal numbers and exponentiation. \emph{Fundamenta Mathematicae} 167 (2001), 173--188.
\href{https://doi.org/10.4064/fm167-2-3}{doi:10.4064/fm167-2-3}.
Erratum, \emph{Fundamenta Mathematicae} 168 (2001), 295--297. Exact ordinal closure estimates must be read with the correction.
\bibitem{Neumann}
B. H. Neumann. On ordered division rings. \emph{Transactions of the American Mathematical Society} 66 (1949), 202--252.
\bibitem{Poonen}
Bjorn Poonen. Maximally complete fields. \emph{L'Enseignement Math\'ematique} (2) 39 (1993), 87--106.
\href{https://math.mit.edu/~poonen/papers/amsval.pdf}{Author's full text}; Corollary 4 supplies the algebraic-closure statement used here.
\bibitem{Ehrlich}
Philip Ehrlich. Number systems with simplicity hierarchies: a generalization of Conway's theory of surreal numbers.
\emph{Journal of Symbolic Logic} 66 (2001), 1231--1258. Corrigendum, 70 (2005), 1022,
\href{https://doi.org/10.2178/jsl/1122038926}{doi:10.2178/jsl/1122038926}.
\bibitem{EKHierarchy}
Philip Ehrlich and Elliot Kaplan. Number systems with simplicity hierarchies II.
\href{https://arxiv.org/abs/1512.04001}{arXiv:1512.04001}; published as \emph{Number Systems with Simplicity Hierarchies: A Generalization of Conway's Theory of Surreal Numbers II}, \emph{Journal of Symbolic Logic} (2018).
\bibitem{EKExp}
Philip Ehrlich and Elliot Kaplan. Surreal ordered exponential fields.
\emph{Journal of Symbolic Logic}, published online 2021.
\href{https://arxiv.org/abs/2002.07739}{arXiv:2002.07739}, especially \S11 and Theorem 11.1 on the canonical surcomplex exponential with omnific-integer kernel.
\bibitem{Williams}
Kameryn J. Williams. Minimum models of second-order set theories.
\emph{Journal of Symbolic Logic} 84 (2019), 589--620.
\href{https://doi.org/10.1017/jsl.2019.27}{doi:10.1017/jsl.2019.27};
\href{https://arxiv.org/abs/1709.03955}{arXiv:1709.03955}. See in particular the definitions of GB, GBC, KM, and ETR.
\bibitem{HamkinsRecursion}
Joel David Hamkins. Determinacy for proper-class clopen games is equivalent to transfinite recursion along proper-class well-founded relations. Author's mathematical exposition, 2014.
\href{https://jdh.hamkins.org/determinacy-for-proper-class-clopen-games-is-equivalent-to-transfinite-recursion-along-proper-class-well-founded-relations/}{Online exposition}. Used for the distinction between set-like recursion and unrestricted class recursion.
\bibitem{ClassForcing}
Victoria Gitman, Joel David Hamkins, Peter Holy, Philipp Schlicht, and Kameryn J. Williams. The exact strength of the class forcing theorem.
\emph{Journal of Symbolic Logic}, published online July 29, 2020.
\href{https://www.cambridge.org/core/journals/journal-of-symbolic-logic/article/abs/exact-strength-of-the-class-forcing-theorem/2295032404B1308A9D4D91065F514FA1}{Publisher's record}. Establishes the role of $\ETR_{\On}$ for the class forcing theorem.
\bibitem{GHJ}
Victoria Gitman, Joel David Hamkins, and Thomas A. Johnstone. What is the theory ZFC without power set?
\href{https://arxiv.org/abs/1110.2430}{arXiv:1110.2430}. Used only for the warning about weakening set-existence and Collection principles.
\bibitem{PakKaliszyk}
Karol P\k{a}k and Cezary Kaliszyk. Conway Normal Form: Bridging Approaches for Comprehensive Formalization of Surreal Numbers.
\href{https://arxiv.org/abs/2410.00065}{arXiv:2410.00065}, 2024. Primary report on the Mizar construction, arithmetic, and normal-form formalization. The carrier $\No$ is real closed, not algebraically closed; the latter property belongs to $\No[i]$.
\bibitem{GambinoAczel}
Nicola Gambino and Peter Aczel. The generalised type-theoretic interpretation of constructive set theory.
\emph{Journal of Symbolic Logic} 71 (2006), 67--103.
\href{https://doi.org/10.2178/jsl/1140641163}{doi:10.2178/jsl/1140641163}.
\bibitem{HoTT}
The Univalent Foundations Program. \emph{Homotopy Type Theory: Univalent Foundations of Mathematics}. Institute for Advanced Study, 2013, \S11.6.
\href{https://homotopytypetheory.org/book/}{Book and source}. The surreal construction uses higher inductive-inductive definitions.
\bibitem{Shulman}
Mike Shulman. The surreals contain the plump ordinals. February 22, 2014.
\href{https://homotopytypetheory.org/2014/02/22/surreals-plump-ordinals/}{Author's exposition on the HoTT website}.
\bibitem{Joseph}
Jean S. Joseph. The Surreal Numbers as an Ordered Commutative Ring with an Apartness: A Development in Univalent Foundations.
\href{https://arxiv.org/abs/1812.00051}{arXiv:1812.00051}, 2018.
\bibitem{RangelMariano}
Dimi Rocha Rangel and Hugo Luiz Mariano. An algebraic (set) theory of surreal numbers, I.
\href{https://arxiv.org/abs/1911.12726}{arXiv:1911.12726}, 2019.
\bibitem{SurrealArithmetic}
Joel David Hamkins. Surreal arithmetic is bi-interpretable with set theory. CUNY Logic Workshop announcement, March 13, 2026.
\href{https://jdh.hamkins.org/surreal-arithmetic-cuny-logic-workshop-march-2026/}{Author's announcement}. Joint work in progress with Junhong Chen and Ruizhi Yang; reported here as an announcement, not used as a proved dependency.
\bibitem{LeanUniverses}
The Lean development team. \emph{The Lean Language Reference}, ``Universes.''
\href{https://lean-lang.org/doc/reference/latest/The-Type-System/Universes/}{Official reference}. Consulted September 21, 2026. Predicativity, noncumulativity, universe polymorphism, and explicit lifting.
\bibitem{LeanAxioms}
The Lean development team. \emph{The Lean Language Reference}, ``Axioms'' and ``Propositions.''
\href{https://lean-lang.org/doc/reference/latest/Axioms/}{Axioms};
\href{https://lean-lang.org/doc/reference/latest/The-Type-System/Propositions/}{Propositions}. Consulted September 21, 2026. Proof irrelevance, classical axioms, \texttt{sorryAx}, and axiom reporting.
\bibitem{LeanGames}
Contributors to \texttt{vihdzp/combinatorial-games}. Lean 4 combinatorial-games and surreal-number library.
\href{https://github.com/vihdzp/combinatorial-games/tree/02b4a908ea2ecfefffecb438f691951a814a5264}{Pinned revision, September 19, 2026}.
Inspected files: \texttt{Surreal/Basic.lean}, \texttt{Surreal/Division.lean}, \texttt{Surreal/HahnSeries/Basic.lean} under \texttt{CombinatorialGames/}, plus \texttt{lean-toolchain} and \texttt{lake-manifest.json}. Source inspection on September 21, 2026; no local build claimed.
\bibitem{MathlibHahn}
Mathlib contributors. \texttt{Mathlib.RingTheory.HahnSeries.Summable}.
\href{https://leanprover-community.github.io/mathlib4_docs/Mathlib/RingTheory/HahnSeries/Summable.html}{Generated primary API documentation}. Consulted September 21, 2026. Formal summable families and \texttt{hsum}; documentation inspection is distinct from compilation against the pinned dependency.
\bibitem{MathlibRealClosed}
Mathlib contributors. \texttt{Mathlib.FieldTheory.IsRealClosed.Basic}.
\href{https://leanprover-community.github.io/mathlib4_docs/Mathlib/FieldTheory/IsRealClosed/Basic.html}{Generated primary API documentation}. Consulted September 21, 2026. Presence of the interface is not a verification of a surreal instance.
\bibitem{MathlibZFC}
Mathlib contributors. \texttt{Mathlib.SetTheory.ZFC.Basic}.
\href{https://leanprover-community.github.io/mathlib4_docs/Mathlib/SetTheory/ZFC/Basic.html}{Generated primary API documentation}. Consulted September 21, 2026. Universe-indexed pre-set quotient, internal set operations, and the technical \texttt{Definable} interface.
\end{thebibliography}
\endgroup
\end{document}
